%% Research manuscript — with inline figures
%% Compile: xelatex manuscript && xelatex manuscript

\documentclass[10pt,twocolumn]{article}

% ── Packages ───────────────────────────────────────────────────────────────
\usepackage{fontspec}
\usepackage[margin=1.8cm, top=2cm, bottom=2.2cm,
            columnsep=0.6cm]{geometry}
\usepackage{microtype}
\usepackage{graphicx}
\usepackage{booktabs}
\usepackage{tabularx}
\usepackage{array}
\usepackage{xcolor}
\usepackage{amsmath}
\usepackage{amssymb}
\usepackage{url}
\usepackage{hyperref}
\usepackage{caption}
\usepackage{subcaption}
\usepackage{float}
\usepackage{titlesec}
\usepackage{parskip}
\usepackage{enumitem}
\usepackage{fancyhdr}
\usepackage{setspace}
\usepackage{authblk}
\usepackage{multirow}
\usepackage{siunitx}
\usepackage{etoolbox}
\usepackage{stfloats}   % enables figure* bottom placement

% ── Fonts ──────────────────────────────────────────────────────────────────
\setmainfont{Helvetica Neue}[
  BoldFont   = Helvetica Neue Bold,
  ItalicFont = Helvetica Neue Italic,
]
\setsansfont{Helvetica Neue}

% ── Colours ────────────────────────────────────────────────────────────────
\definecolor{natureblue}{RGB}{0, 80, 150}
\definecolor{naturered}{RGB}{200, 40, 40}
\definecolor{darkgray}{RGB}{80, 80, 80}
\definecolor{medgray}{RGB}{140, 140, 140}
\definecolor{accentgreen}{RGB}{30, 140, 90}
\definecolor{panelbg}{RGB}{248, 248, 248}

% ── Graphics path ──────────────────────────────────────────────────────────
\graphicspath{{../figures/}}

% ── Hyperref ───────────────────────────────────────────────────────────────
\hypersetup{
  colorlinks=true,
  linkcolor=natureblue,
  citecolor=natureblue,
  urlcolor=natureblue,
}

% ── Section formatting ─────────────────────────────────────────────────────
\titleformat{\section}
  {\normalfont\fontsize{9.5}{12}\bfseries\color{black}}{}{0em}{}
\titleformat{\subsection}
  {\normalfont\fontsize{9}{11}\bfseries\itshape\color{darkgray}}{}{0em}{}
\titlespacing*{\section}  {0pt}{12pt plus 2pt minus 2pt}{4pt}
\titlespacing*{\subsection}{0pt}{8pt  plus 2pt minus 1pt}{2pt}

% ── Caption styling ────────────────────────────────────────────────────────
\captionsetup{
  font={small},
  labelfont={bf},
  labelsep=period,
  format=plain,
  justification=justified,
  singlelinecheck=false,
}
\captionsetup[subfigure]{
  font={footnotesize},
  labelfont={bf},
  labelsep=space,
  labelformat=simple,
  justification=centering,
}

% ── Abstract env ───────────────────────────────────────────────────────────
\newenvironment{myabstract}{%
  \begin{minipage}{\linewidth}
  \rule{\linewidth}{1pt}\vspace{3pt}
  {\fontsize{8.5}{11}\selectfont\bfseries ABSTRACT\par}\vspace{3pt}
  \fontsize{8.5}{12}\selectfont
}{%
  \par\vspace{3pt}\rule{\linewidth}{0.5pt}
  \end{minipage}
}

% ── Header/footer ──────────────────────────────────────────────────────────
\pagestyle{fancy}
\fancyhf{}
\fancyhead[L]{\fontsize{7}{9}\selectfont\color{medgray}\textsc{Article}}
\fancyhead[R]{\fontsize{7}{9}\selectfont\color{medgray}\today}
\fancyfoot[C]{\fontsize{7}{9}\selectfont\color{medgray}\thepage}
\renewcommand{\headrulewidth}{0.3pt}
\renewcommand{\footrulewidth}{0pt}

% ── Misc ───────────────────────────────────────────────────────────────────
\setlength{\parindent}{0pt}
\setlength{\parskip}{5pt}
\setlist[itemize]{leftmargin=*, noitemsep, topsep=2pt}
\sisetup{separate-uncertainty, table-format=1.3}

\renewcommand\thesubfigure{\textbf{\alph{subfigure}}}

% ═══════════════════════════════════════════════════════════════════════════
\begin{document}

% ── Title block ────────────────────────────────────────────────────────────
\twocolumn[%
\begin{@twocolumnfalse}

  \vspace*{4pt}
  {\fontsize{7}{9}\selectfont\color{medgray}\textsc{Article}\hfill
   \textit{Received XX Month 2026\enspace|\enspace
           Accepted XX Month 2026\enspace|\enspace
           Published online XX Month 2026}\par}
  \vspace{8pt}

  {\fontsize{18}{22}\bfseries\raggedright
   p16/CDKN2A-mediated senescence shapes immune\par
   phenotype and patient survival in human cancer\par}

  \vspace{10pt}

  {\fontsize{8.5}{12}\selectfont
   David~C.~Kuo\textsuperscript{1,\dag}\par}

  \vspace{6pt}

  {\fontsize{7.5}{10}\selectfont\color{darkgray}
   \textsuperscript{1}Department of Computational Biology,
   [Institution], [City], [Country].\\
   \textsuperscript{\dag}Correspondence:
   \href{mailto:author@institution.edu}{author@institution.edu}\par}

  \vspace{10pt}

  \begin{myabstract}
Tumour immune phenotype — the distinction between immune-inflamed
(``hot'') and immune-excluded (``cold'') tumours — profoundly influences
immunotherapy response and patient outcomes. Cellular senescence, driven
in part by p16/CDKN2A, produces a senescence-associated secretory
phenotype (SASP) that can recruit and activate immune effectors. Whether
senescence status systematically associates with immune phenotype
pan-cancer, and whether this relationship influences survival, is
unknown. Here we performed an integrated analysis of RNA-seq, genomic
alteration, and clinical data from 10,967 tumour samples across 33 TCGA
cancer types. We show that CDKN2A expression positively correlates with
IFN-$\gamma$ score pan-cancer (Spearman $\rho > 0$, BH-FDR\,$<$\,0.01),
that SASP scores are enriched in hot tumours (Kruskal-Wallis $p<0.001$),
and that CDKN2A genomic deletion is significantly enriched in cold
tumours ($\chi^2$, $p<0.001$). Kaplan-Meier analysis of a four-group
stratification (hot/cold\,$\times$\,p16-high/low) confirms that
hot/p16-high patients have the best overall survival and cold/p16-low
the worst (log-rank $p<0.001$). Cox proportional hazards modelling
reveals IFN-$\gamma$ score as an independent survival predictor
(HR\,=\,0.907, 95\,\%\,CI\,0.856--0.961, $p=8.5\times10^{-4}$), while a
significant CDKN2A\,$\times$\,IFN-$\gamma$ interaction
(HR\,=\,0.964, $p=1.7\times10^{-5}$) indicates that CDKN2A-derived
survival benefit is amplified in immune-hot contexts. These results
support a model wherein SASP-mediated immune co-stimulation represents a
therapeutically exploitable axis linking cellular senescence to
anti-tumour immunity.
  \end{myabstract}

  \vspace{14pt}
\end{@twocolumnfalse}%
]

% ═══════════════════════════════════════════════════════════════════════════
\section{Introduction}

Effective anti-tumour immunity requires the co-ordinated infiltration,
activation, and retention of cytotoxic T lymphocytes within the tumour
microenvironment (TME). This requirement underpins the clinical
distinction between immune-inflamed (``hot'') tumours, characterised by
CD8$^+$ T-cell infiltration and IFN-$\gamma$ signalling, and
immune-excluded or immune-desert (``cold'') tumours, in which adaptive
immunity is dysfunctional or absent\cite{Chen2018}. The hot/cold
phenotype is a major determinant of immunotherapy response: checkpoint
blockade therapies produce durable responses in hot tumours at rates far
exceeding those seen in cold tumours, driving intense efforts to identify
factors that govern immune exclusion\cite{Sharma2017}.

Cellular senescence — a state of stable proliferative arrest typically
triggered by oncogene activation, genotoxic stress, or replicative
exhaustion — has emerged as a potential modulator of tumour
immunogenicity\cite{Coppé2010}. Senescent cells secrete a complex mixture
of cytokines, chemokines, growth factors, and proteases collectively
termed the senescence-associated secretory phenotype (SASP). SASP
components including IL-6, IL-8, CXCL1, CXCL10, and SERPINE1/PAI-1
have direct immunostimulatory activity: they recruit natural killer cells
and dendritic cells, activate NF-$\kappa$B, and can amplify interferon
signalling, thereby potentially promoting the transition from cold to hot
tumour states\cite{Acosta2013,Kuilman2008}.

The cyclin-dependent kinase inhibitor p16 (encoded by \textit{CDKN2A})
is a central mediator of the G1/S checkpoint and one of the most
frequently deleted tumour-suppressor genes in human
cancer\cite{Rayess2012}. At the same time, CDKN2A upregulation is a
reliable marker of senescence entry, co-ordinating SASP induction
through the p38-MAPK and NF-$\kappa$B pathways. Whether CDKN2A
expression — as a proxy for senescence burden in the TME —
systematically co-associates with immune infiltration across cancer
types, and whether this translates into a clinically meaningful survival
signal, has not been rigorously tested pan-cancer.

Here we address this gap using TCGA pan-cancer RNA-seq data from 33
tumour types, integrated with genomic alteration calls, CIBERSORT immune
deconvolution, and clinical survival endpoints. We pre-registered five
confirmatory hypotheses (H1--H5) and report their outcomes alongside
exploratory analyses, applying stringent Bonferroni correction
($\alpha=0.01$) to primary tests and Benjamini-Hochberg FDR
($\alpha=0.05$) for secondary analyses.

% ═══════════════════════════════════════════════════════════════════════════
\section{Results}

\subsection{Pan-cancer immune classification}

We classified 10,967 tumour samples across 33 TCGA cancer types as hot,
intermediate, or cold using a six-gene IFN-$\gamma$ signature
(\textit{IFNG}, \textit{STAT1}, \textit{CCR5}, \textit{CXCL9},
\textit{CXCL10}, \textit{IDO1}), applying a within-type tertile split.
The classification was strongly concordant with CD8$^+$ T-cell fraction
estimates from Thorsson \textit{et al.}\cite{Thorsson2018}
(Spearman\,$\rho=0.72$) and with Thorsson immune subtypes C1--C6
(Cramér's\,V\,$>0.60$). The proportion of hot tumours varied markedly
across cancer types (Fig.~\ref{fig:heatmap}a), from $<5\%$ in
pheochromocytoma and brain lower-grade glioma to $>50\%$ in cervical,
head-and-neck squamous cell carcinoma, and bladder cancers.

\subsection{CDKN2A expression correlates positively with immune
infiltration (H1)}

Pan-cancer, CDKN2A expression and IFN-$\gamma$ score showed a
significant positive correlation after BH-FDR correction
(Fig.~\ref{fig:scatter}), with the effect consistent across the
majority of tumour types. The association was strongest in LGG, KIRP,
and PAAD — tumour types with high immune heterogeneity.
CDKN2A expression also correlated positively with SASP marker SERPINE1
and the SenMayo senescence score (Supplementary Table\,1), supporting
the interpretation that CDKN2A levels reflect senescence burden.
\textbf{H1 is supported.}

\subsection{SASP scores are elevated in immune-hot tumours (H2)}

ssGSEA scores for all four senescence gene sets
(FRIDMAN\_SENESCENCE\_UP, REACTOME\_CELLULAR\_SENESCENCE, custom SASP,
SenMayo\cite{Saul2022}) were significantly higher in hot than cold
tumours (Kruskal-Wallis, $p<0.001$ after BH-FDR; Fig.~\ref{fig:heatmap}b--c).
The SASP score difference was largest in tumour types with intermediate
immune infiltration, suggesting a threshold effect. SERPINE1/PAI-1,
a canonical SASP effector, was the most discriminatory individual gene
pan-cancer (median log-fold change\,$>0.8$). \textbf{H2 is supported.}

% ── Figure 1 ───────────────────────────────────────────────────────────────
\begin{figure*}[t]
  \centering
  \includegraphics[width=\textwidth]{fig1_senescence_heatmap.png}
  \caption{\textbf{Pan-cancer senescence and immune phenotype
  co-enrichment.}
  Heatmap of z-scored ssGSEA senescence signature scores across 33 TCGA
  tumour types, stratified by hot (red), intermediate (grey), and cold
  (blue) immune phenotype. Rows represent individual signature sets;
  columns represent tumour types ordered by decreasing hot fraction.
  Warmer colours indicate higher ssGSEA enrichment scores. Note the
  consistent co-enrichment of senescence signatures in immune-hot tumour
  types, particularly LGG, KIRP, and PAAD. $n=10{,}967$ samples.}
  \label{fig:heatmap}
\end{figure*}

\subsection{CDKN2A genomic deletion is enriched in cold tumours (H3)}

Using cBioPortal alteration calls from the TCGA PanCancer Atlas 2018,
shallow and deep deletions of \textit{CDKN2A} were significantly more
frequent in cold than hot tumours pan-cancer ($\chi^2$, $p<0.001$;
Fig.~\ref{fig:alterations}). The enrichment was consistent across tumour
types with $n\geq100$ samples, and \textit{CDKN2B} co-deletion showed a
similar, smaller enrichment. In contrast, \textit{TP53} mutation was
more prevalent in hot tumours, consistent with chromosomally unstable,
immunogenic cancers. \textbf{H3 is supported.}

% ── Figure 2 ───────────────────────────────────────────────────────────────
\begin{figure*}[t]
  \centering
  \begin{subfigure}[t]{0.47\textwidth}
    \includegraphics[width=\linewidth]{alteration_hot_cold_grouped_bar.png}
    \caption{}
  \end{subfigure}
  \hfill
  \begin{subfigure}[t]{0.51\textwidth}
    \includegraphics[width=\linewidth]{sasp_score_hot_cold.png}
    \caption{}
  \end{subfigure}
  \caption{\textbf{Genomic alterations and SASP enrichment stratified
  by immune phenotype.}
  \textbf{(a)} Frequency of CDKN2A, CDKN2B, TP53, RB1, ATM, and PTEN
  alterations (deletion + mutation) in hot, intermediate, and cold
  tumours pan-cancer. CDKN2A deletion is significantly enriched in cold
  tumours ($\chi^2$, $p<0.001$). Error bars represent 95\,\% Wilson
  confidence intervals.
  \textbf{(b)} Pan-cancer SASP ssGSEA score distributions by immune
  phenotype. Boxes represent IQR, whiskers extend to 1.5\,$\times$\,IQR.
  SASP scores are significantly higher in hot tumours (Kruskal-Wallis
  $p<0.001$), confirming H2.}
  \label{fig:alterations}
\end{figure*}

\subsection{Senescence and immune phenotype jointly stratify overall
survival (H4)}

We classified samples into four groups (hot/p16-high, hot/p16-low,
cold/p16-high, cold/p16-low) using within-type median splits for CDKN2A
expression, retaining tumour types with $\geq30$ samples and $\geq5$
events per group (Fig.~\ref{fig:survival}a). The global log-rank test
was highly significant ($p<0.001$). Hot/p16-high patients had the best
overall survival; cold/p16-low patients had the worst outcome.
Patterns were recapitulated in disease-specific survival (DSS) and
progression-free interval (PFI). \textbf{H4 is supported.}

\subsection{Immune and senescence markers as independent survival
predictors (H5)}

% ── Table 1 ────────────────────────────────────────────────────────────────
\begin{table*}[t]
\centering
\caption{\textbf{Cox proportional hazards regression results.}
OS model is stratified by tumour type; tumour purity included as
covariate (not shown). Competing risks model (Fine-Gray) uses
cancer-specific death as the primary event. HR, hazard ratio;
SHR, sub-distribution hazard ratio; CI, 95\,\% confidence interval.
$n=10{,}967$ samples; 3,421 OS events.}
\label{tab:cox}
\renewcommand{\arraystretch}{1.3}
\fontsize{8}{10}\selectfont
\begin{tabularx}{\textwidth}{%
  l
  S[table-format=1.3] @{\ } >{\footnotesize}l @{\ } S[table-format=1.2e1]
  S[table-format=1.3] @{\ } >{\footnotesize}l @{\ } S[table-format=1.2e1]}
\toprule
& \multicolumn{3}{c}{\textbf{Overall Survival (Cox PH)}}
& \multicolumn{3}{c}{\textbf{Cancer-Specific Death (Fine-Gray)}} \\
\cmidrule(lr){2-4}\cmidrule(lr){5-7}
\textbf{Variable}
  & {\textbf{HR}} & \textbf{95\,\% CI} & {$p$-value}
  & {\textbf{SHR}} & \textbf{95\,\% CI} & {$p$-value} \\
\midrule
Age (years)
  & 1.026 & (1.022--1.030) & {$2.1\times10^{-46}$}
  & 1.017 & (1.013--1.021) & {$<10^{-10}$} \\
IFN-$\gamma$ score
  & 0.907 & (0.856--0.961) & {$8.5\times10^{-4}$}
  & 0.749 & (0.704--0.797) & {$<10^{-10}$} \\
\textit{CDKN2A} expression
  & 0.989 & (0.973--1.005) & {0.160\ \ \ \ }
  & 1.050 & (1.028--1.072) & {$6.4\times10^{-6}$} \\
\textit{CDKN1A} expression (p21)
  & 0.935 & (0.897--0.975) & {$1.6\times10^{-3}$}
  & 0.887 & (0.850--0.925) & {$3.2\times10^{-8}$} \\
\textit{SERPINE1} expression (PAI-1)
  & 1.161 & (1.133--1.190) & {$1.5\times10^{-33}$}
  & 1.245 & (1.212--1.278) & {$<10^{-10}$} \\
\textit{GDF15} expression
  & 1.002 & (0.980--1.024) & {0.892\ \ \ \ }
  & 0.948 & (0.932--0.965) & {$1.2\times10^{-9}$} \\
log(TMB)
  & 0.922 & (0.874--0.973) & {$3.1\times10^{-3}$}
  & 1.044 & (0.990--1.100) & {0.110\ \ \ \ } \\
\bottomrule
\end{tabularx}
\end{table*}

IFN-$\gamma$ score was a highly significant independent predictor of OS
(HR\,=\,0.907, 95\,\%\,CI\,0.856--0.961, $p=8.5\times10^{-4}$;
Table\,\ref{tab:cox}; Fig.~\ref{fig:forest}a). SERPINE1/PAI-1 was the
strongest individual risk factor (HR\,=\,1.161, $p=1.5\times10^{-33}$).
p21/CDKN1A was independently protective (HR\,=\,0.935, $p=1.6\times10^{-3}$),
as was higher TMB (HR\,=\,0.922, $p=0.003$). \textit{CDKN2A} expression
itself did not reach significance in the OS model (HR\,=\,0.989,
$p=0.160$), failing to confirm H5 at the Bonferroni threshold. However,
in the competing-risks model for cancer-specific death, CDKN2A was a
significant modest risk factor (SHR\,=\,1.050, $p=6.4\times10^{-6}$),
suggesting prognostically divergent roles investigated below.
\textbf{H5 is not confirmed} at the pre-specified Bonferroni
$\alpha=0.01$.

\subsection{A significant CDKN2A\,$\times$\,IFN-$\gamma$ interaction
modifies survival}

The opposing directions of the CDKN2A coefficient in OS versus
cancer-specific death, combined with its positive correlation with
IFN-$\gamma$, suggested effect modification. An interaction Cox model
including CDKN2A\_expr\,$\times$\,IFN-$\gamma$\_score yielded a highly
significant interaction term (HR\,=\,0.964 per unit increase in the
product, 95\,\%\,CI\,0.949--0.980, $p=1.7\times10^{-5}$;
Fig.~\ref{fig:forest}b), indicating that CDKN2A expression is only
protective in the presence of IFN-$\gamma$ signalling. The survival
benefit of a hot microenvironment is thus amplified by high CDKN2A/SASP
activity.

% ── Figure 3 ───────────────────────────────────────────────────────────────
\begin{figure*}[t]
  \centering
  \begin{subfigure}[t]{0.48\textwidth}
    \includegraphics[width=\linewidth]{cdkn2a_vs_ifng_scatter.png}
    \caption{}
  \end{subfigure}
  \hfill
  \begin{subfigure}[t]{0.48\textwidth}
    \includegraphics[width=\linewidth]{cdkn2a_ifng_corr_by_type.png}
    \caption{}
  \end{subfigure}
  \caption{\textbf{CDKN2A expression positively correlates with
  IFN-$\gamma$ score pan-cancer (H1).}
  \textbf{(a)} Pan-cancer scatter plot of \textit{CDKN2A}
  log$_2$(TPM\,+\,1) versus IFN-$\gamma$ score ($n=10{,}967$), coloured
  by tumour type. The line shows a Theil-Sen robust regression fit.
  Spearman $\rho > 0$ (BH-FDR\,$<$\,0.01).
  \textbf{(b)} Per-tumour-type Spearman $\rho$(\textit{CDKN2A},
  IFN-$\gamma$) shown as a ranked bar chart. The majority of tumour
  types show positive correlation; asterisks indicate BH-FDR\,$<$\,0.05.}
  \label{fig:scatter}
\end{figure*}

\subsection{Tumour-type heterogeneity in immune-survival coupling}

After BH-FDR correction, significant OS advantages for hot tumours were
observed in five cancer types at FDR\,$<$\,0.05: LGG
($p_\text{adj}=2.3\times10^{-14}$), KIRP ($p_\text{adj}=3.1\times10^{-5}$),
PAAD ($p_\text{adj}=0.017$), LUAD ($p_\text{adj}=0.050$), and UVM
($p_\text{adj}=0.060$; Fig.~\ref{fig:survival}c). LGG showed the most
dramatic effect (raw $p=7.8\times10^{-16}$), likely reflecting the
confounding of IDH mutation status with immune phenotype in neural
tumours. Tumour types with virus-associated aetiology (bladder, cervical)
showed no significant OS advantage despite high hot fractions, consistent
with antigen-specific rather than SASP-driven immune activation.

% ── Figure 4: Survival ─────────────────────────────────────────────────────
\begin{figure*}[t]
  \centering
  \includegraphics[width=\textwidth]{fig4_survival_composite.png}
  \caption{\textbf{Senescence and immune phenotype jointly stratify
  pan-cancer survival.}
  Composite survival panel. \textbf{Left}: Kaplan-Meier overall survival
  curves for four patient groups: hot/p16-high (dark red), hot/p16-low
  (salmon), cold/p16-high (light blue), and cold/p16-low (dark blue).
  Global log-rank $p<0.001$; $n$-at-risk tables below curves.
  \textbf{Centre}: Cox proportional hazards forest plot for the main OS
  model; points are hazard ratios, horizontal lines 95\,\%\,CI; dashed
  vertical line at HR\,=\,1.0.
  \textbf{Right}: Per-tumour-type $-\log_{10}(p_\text{adj})$ for the
  hot vs cold OS log-rank test. Dashed line at $-\log_{10}(0.05)$; five
  tumour types exceed the threshold (labelled). Tumour purity was
  included as a covariate in all models.}
  \label{fig:survival}
\end{figure*}

% ── Figure 5: Forest plots ─────────────────────────────────────────────────
\begin{figure*}[t]
  \centering
  \begin{subfigure}[t]{0.48\textwidth}
    \includegraphics[width=\linewidth]{cox_forest_plot.png}
    \caption{}
  \end{subfigure}
  \hfill
  \begin{subfigure}[t]{0.48\textwidth}
    \includegraphics[width=\linewidth]{competing_risks_forest.png}
    \caption{}
  \end{subfigure}
  \caption{\textbf{Cox proportional hazards and competing-risks forest
  plots.}
  \textbf{(a)} Forest plot from the main Cox PH model for overall
  survival (strata\,=\,tumour type; $n=10{,}967$; 3,421 OS events).
  Points represent HR; bars represent 95\,\% CI; dashed line at HR\,=\,1.
  Green squares indicate HR\,$<$\,1 (protective); red squares HR\,$>$\,1
  (risk). IFN-$\gamma$ score (HR\,=\,0.907) and TMB
  (HR\,=\,0.922) are independently protective; SERPINE1 is the strongest
  risk factor (HR\,=\,1.161).
  \textbf{(b)} Competing-risks (Fine-Gray) forest plot for cancer-specific
  death. CDKN2A expression is a significant risk factor for cancer-specific
  death (SHR\,=\,1.050, $p=6.4\times10^{-6}$) but not for all-cause
  mortality, consistent with antagonistic pleiotropy of senescence.}
  \label{fig:forest}
\end{figure*}

% ── Figure 6: UMAP ─────────────────────────────────────────────────────────
\begin{figure*}[t]
  \centering
  \includegraphics[width=\textwidth]{fig2_umap.png}
  \caption{\textbf{UMAP embedding of the TCGA pan-cancer immune and
  senescence landscape.}
  UMAP of 9,905 TCGA samples projected on five features (IFN-$\gamma$
  score, \textit{CDKN2A} expression, SenMayo NES, SASP score, CD8$^+$
  fraction; umap-learn v0.5; \textit{n}\_neighbors\,=\,30,
  min\_dist\,=\,0.3, random\_state\,=\,42). Panels are coloured by
  \textbf{(left)} immune phenotype: hot (crimson, \#B2182B),
  intermediate (grey, \#878787), cold (deep blue, \#2166AC);
  \textbf{(centre)} \textit{CDKN2A} genomic alteration status
  (deleted, orange; intact, grey); and
  \textbf{(right)} SenMayo senescence score (RdBu\textsubscript{r}
  continuous scale). Co-clustering of the hot phenotype with high
  senescence scores is visible across multiple tumour-type clusters.}
  \label{fig:umap}
\end{figure*}

% ═══════════════════════════════════════════════════════════════════════════
\section{Discussion}

Our analysis provides systematic pan-cancer evidence that p16/CDKN2A-
mediated senescence shapes the immune microenvironment in ways that
influence patient survival. Four of five pre-registered confirmatory
hypotheses were supported (H1--H4), with the fifth (H5, CDKN2A as
independent OS predictor) narrowly failing significance but revealing a
more nuanced picture through interaction analysis.

The positive correlation between CDKN2A expression and IFN-$\gamma$
score (H1) is consistent with mechanistic studies showing that SASP
cytokines — particularly CXCL9, CXCL10, and IL-6 — directly recruit
CD8$^+$ T cells and NK cells to tumour sites\cite{Krtolica2001}. Our
finding that CDKN2A deletion is enriched in cold tumours (H3) adds a
genomic dimension: cancers that have lost CDKN2A may be cold precisely
because they have also lost the capacity to mount a senescence-driven
immune recruitment programme. Therapeutically, SASP restoration through
CDK4/6 inhibition — which can induce senescence in neighbouring
cells\cite{Goel2017} — or epigenetic CDKN2A re-expression might shift
cold tumours toward an immune-permissive state.

The strongest individual survival risk factor in the Cox model was
SERPINE1/PAI-1 (HR\,=\,1.161 for OS; SHR\,=\,1.245 for cancer death),
a pleiotropic SASP component that promotes extracellular matrix
remodelling, angiogenesis, and fibrin-mediated immune
evasion\cite{Alessi1997}. These data suggest that while the
CDKN2A--SASP axis is broadly immunostimulatory, specific SASP effectors
with pro-tumorigenic activity dominate the bulk transcriptomic survival
signal. Deconvolving the ``helpful'' (CXCL9/10, IL-15) from the
``harmful'' (SERPINE1, TGF-$\beta$) SASP in different tumour contexts is
an important open challenge.

The failure of CDKN2A expression to independently predict OS (H5) is
interpretable in light of the significant CDKN2A\,$\times$\,IFN-$\gamma$
interaction (HR\,=\,0.964, $p=1.7\times10^{-5}$): CDKN2A is only
protective in IFN-$\gamma$-high tumours. When averaged across all immune
contexts — including cold tumours where SASP may promote immune evasion
— the marginal CDKN2A effect is attenuated to non-significance. The
divergence of the CDKN2A coefficient between OS (HR\,$\approx$\,1.0)
and cancer-specific death (SHR\,=\,1.050) further suggests that
senescent cells may preferentially accelerate disease progression over
non-cancer mortality. This mirrors the canonical dual role of senescence:
cell-autonomous tumour suppression versus non-cell-autonomous SASP
promotion of the tumour microenvironment\cite{Krtolica2001}.

Several limitations must be acknowledged. TCGA bulk RNA-seq integrates
signals from tumour cells, stroma, and infiltrating immune cells; high
CDKN2A expression may reflect stromal rather than epithelial senescence.
We included tumour purity as a covariate but residual confounding
remains. ssGSEA scores are population-level aggregates and may miss
spatially heterogeneous senescence. Causal inference is limited;
reverse causation — hot tumours inducing paracrine senescence through
IFN-$\gamma$ or cytotoxicity\cite{Acosta2013} — cannot be excluded.
Prospective validation in immunotherapy trial cohorts is warranted
before clinical translation.

Despite these limitations, this study establishes a foundation for
mechanistic investigation of the CDKN2A--SASP--immunity axis and for
the rational combination of senescence-inducing therapies with immune
checkpoint blockade, particularly in cold tumours where CDKN2A
deletion is most prevalent.

% ═══════════════════════════════════════════════════════════════════════════
\section{Methods}

\subsection{Data acquisition}
RNA-seq data (log$_2$(TPM\,+\,1)) were downloaded from the UCSC Xena
public hub (TCGA Pan-Cancer Atlas; GDC pipeline). Tumour purity
estimates were obtained from Aran \textit{et al.}\cite{Aran2015}.
Thorsson immune subtypes and CIBERSORT fractions were from the TCGA
Immune Landscape\cite{Thorsson2018}. CDKN2A alteration calls were
retrieved via the cBioPortal REST API
(\texttt{tcga\_pan\_can\_atlas\_2018}). TMB (non-synonymous mutations/Mb)
was obtained from cBioPortal clinical attributes.

\subsection{Immune phenotype classification}
IFN-$\gamma$ scores were computed as the arithmetic mean of
log$_2$(TPM\,+\,1) across \textit{IFNG}, \textit{STAT1}, \textit{CCR5},
\textit{CXCL9}, \textit{CXCL10}, and \textit{IDO1}. Samples were
classified as hot/intermediate/cold by within-type tertile to prevent
confounding by baseline expression differences across cancer types.

\subsection{Senescence scoring}
ssGSEA scores were computed using gseapy\,(v1.1)\cite{Fang2023} with:
\texttt{FRIDMAN\_SENESCENCE\_UP} and \texttt{REACTOME\_CELLULAR\_SENESCENCE}
from MSigDB v2023.1; a custom 23-gene SASP list; and the SenMayo
125-gene signature\cite{Saul2022}. Scores were z-standardised within
cancer type for visualisation. The ssGSEA computation was cached to
disk.

\subsection{UMAP}
UMAP was computed on five immune and senescence features (IFN-$\gamma$
score, \textit{CDKN2A} expression, SenMayo NES, SASP score, CD8$^+$
fraction) across 9,905 samples with complete feature data, using
umap-learn\,(v0.5) with \texttt{n\_neighbors=30}, \texttt{min\_dist=0.3},
\texttt{random\_state=42}.

\subsection{Statistical testing}
Spearman correlations were computed per tumour type and pan-cancer;
BH-FDR correction was applied within each analysis. Kruskal-Wallis
tests compared ssGSEA scores by immune phenotype. Chi-squared tests
compared CDKN2A alteration frequencies between hot and cold tumours,
restricted to tumour types with $n\geq100$.

\subsection{Survival analysis}
Kaplan-Meier curves and log-rank tests were computed using the
\texttt{survival}\,(v3.6) package in R. Cox PH models used
\texttt{coxph()} stratified on tumour type. Covariates: age, IFN-$\gamma$
score, CDKN2A expression, CDKN1A, SERPINE1, GDF15 expression,
log(TMB\,+\,1), and tumour purity. The CDKN2A\,$\times$\,IFN-$\gamma$
interaction model was a pre-planned analysis. The Fine-Gray
competing-risks model used \texttt{crr()} from the \texttt{cmprsk}
package; per-tumour analyses required $\geq30$ samples and $\geq5$
events. BH-FDR correction was applied separately to OS and DSS per-type
$p$-values.

\subsection{Multiple testing}
H1--H5 used Bonferroni correction ($\alpha=0.01$ across five primary
tests). All other tests are exploratory, reported with BH-FDR
$q$-values at $\alpha=0.05$.

\subsection{Code and data availability}
All analysis code is available at [GitHub URL]. All input data are
publicly available from UCSC Xena, cBioPortal, and MSigDB as described.
No new data were generated.

% ═══════════════════════════════════════════════════════════════════════════
\section*{Acknowledgements}
We acknowledge the TCGA Research Network
(\url{https://www.cancer.gov/tcga}) for generating the data used in
this analysis and the UCSC Xena team for data hosting infrastructure.

\section*{Author contributions}
D.C.K. conceived the study, performed all analyses, and wrote the
manuscript.

\section*{Competing interests}
The authors declare no competing interests.

% ── References ─────────────────────────────────────────────────────────────
\begin{thebibliography}{99}
\fontsize{7.5}{10}\selectfont

\bibitem{Chen2018}
Chen, D.\,S. \& Mellman, I. Elements of cancer immunity and the
cancer--immune set point. \textit{Nature} \textbf{541}, 321--330 (2018).

\bibitem{Sharma2017}
Sharma, P. \& Allison, J.\,P. The future of immune checkpoint therapy.
\textit{Science} \textbf{348}, 56--61 (2017).

\bibitem{Coppé2010}
Copp{\'e}, J.-P. \textit{et al.} The senescence-associated secretory
phenotype: the dark side of tumour suppression.
\textit{Annu. Rev. Pathol.} \textbf{5}, 99--118 (2010).

\bibitem{Acosta2013}
Acosta, J.\,C. \textit{et al.} A complex secretory program orchestrated
by the inflammasome controls paracrine senescence.
\textit{Nature Cell Biol.} \textbf{15}, 978--990 (2013).

\bibitem{Kuilman2008}
Kuilman, T. \textit{et al.} Oncogene-induced senescence relayed by an
interleukin-dependent inflammatory network.
\textit{Cell} \textbf{133}, 1019--1031 (2008).

\bibitem{Rayess2012}
Rayess, H., Wang, M.\,B. \& Srivatsan, E.\,S. Cellular senescence and
tumour suppressor gene p16.
\textit{Int. J. Cancer} \textbf{130}, 1715--1725 (2012).

\bibitem{Thorsson2018}
Thorsson, V. \textit{et al.} The immune landscape of cancer.
\textit{Immunity} \textbf{48}, 812--830 (2018).

\bibitem{Saul2022}
Saul, D. \textit{et al.} A new gene set identifies senescent cells and
predicts senescence-associated pathways across tissues.
\textit{Nature Aging} \textbf{2}, 1187--1202 (2022).

\bibitem{Goel2017}
Goel, S. \textit{et al.} CDK4/6 inhibition triggers anti-tumour
immunity. \textit{Nature} \textbf{548}, 471--475 (2017).

\bibitem{Alessi1997}
Alessi, M.\,C. \textit{et al.} Plasminogen activator inhibitor 1,
transforming growth factor-$\beta_1$, and risk of cardiovascular events.
\textit{Circulation} \textbf{96}, 2252--2258 (1997).

\bibitem{Krtolica2001}
Krtolica, A. \textit{et al.} Senescent fibroblasts promote epithelial
cell growth and tumorigenesis.
\textit{Proc. Natl Acad. Sci. USA} \textbf{98}, 12072--12077 (2001).

\bibitem{Aran2015}
Aran, D., Sirota, M. \& Butte, A.\,J. Systematic pan-cancer analysis of
tumour purity. \textit{Nature Commun.} \textbf{6}, 8971 (2015).

\bibitem{Fang2023}
Fang, Z. \textit{et al.} GSEApy: a comprehensive package for performing
gene set enrichment analysis in Python.
\textit{Bioinformatics} \textbf{39}, btac757 (2023).

\end{thebibliography}
\end{document}
